\documentclass[12pt]{article}
\usepackage[T1]{fontenc}
\usepackage[utf8]{inputenc}
\usepackage{lmodern}
\usepackage{textcomp}
\usepackage{enumitem}
\usepackage{geometry}
\geometry{margin=1in}
\begin{document}
\begin{center}
Answer Key - History - Undergraduate Level Assessment 24 \
\small (Answers are ordered exactly as in the question paper.)
\end{center}

\section*{Multiple Choice}
\begin{enumerate}[label=\arabic*.]
\item \textbf{Answer:} A
\\ \textit{Options:} A) elite mismanagement; B) resource depletion; C) natural disasters; D) cultural assimilation; E) technological advancement
\\ \textit{Note:} The correct answer is elite mismanagement because Joseph Tainter's analysis of the collapse of ancient civilizations focuses primarily on factors such as diminishing returns on investment in complexity and resource depletion, rather than specifically attributing collapse to the mismanagement by elite groups.
\item \textbf{Answer:} A
\\ \textit{Options:} A) 600,000 to 300,000 years ago; B) 100,000 years ago to present; C) 400,000 years ago to present; D) 800,000 to 10,000 years ago; E) 500,000 to 200,000 years ago
\\ \textit{Note:} The correct answer is supported by fossil and genetic evidence indicating that anatomically modern Homo sapiens emerged in Africa between 600,000 and 300,000 years ago, marking a significant period in human evolution.
\item \textbf{Answer:} A
\\ \textit{Options:} A) 40,000-45,000 years ago; B) 60,000-65,000 years ago; C) 70,000-75,000 years ago; D) 15,000-20,000 years ago; E) 10,000-12,000 years ago
\\ \textit{Note:} The correct answer of 40,000-45,000 years ago is supported by archaeological evidence indicating that Aboriginal peoples have inhabited the Australian continent for at least this duration, demonstrating their adaptation to the harsh interior environments long before European contact.
\item \textbf{Answer:} B
\\ \textit{Options:} A) the discovery of figurines made from clay at the Trinil site.; B) recovery of a shell at the Trinil site with an intentionally engraved series of lines.; C) the use of complex tools requiring a high level of cognitive function.; D) cave paintings and elaborate burials with numerous grave goods at the Trinil site.; E) all of the above.
\\ \textit{Note:} The correct answer is supported by the fact that the recovery of an intentionally engraved shell at the Trinil site indicates that Homo erectus may have engaged in symbolic expression, suggesting the presence of an aesthetic sense and the potential for artistic creation.
\item \textbf{Answer:} A
\\ \textit{Options:} A) crop irrigation systems; B) the creation of fishing nets; C) European use of bow and arrow; D) domestication of animals; E) canoes
\\ \textit{Note:} Crop irrigation systems are not an innovation of the Mesolithic period, as they emerged later during the Neolithic period when more advanced agricultural practices were developed.
\end{enumerate}

\section*{True / False}
\begin{enumerate}[label=\arabic*.]
\setcounter{enumi}{5}
\item \textbf{Answer:} False
\\ \textit{Options:} A) True; B) False
\\ \textit{Note:} The correct answer is: the air inside is heated and expands causing an upward buoyancy force
\item \textbf{Answer:} False
\\ \textit{Options:} A) True; B) False
\\ \textit{Note:} The correct answer is: crack
\item \textbf{Answer:} True
\\ \textit{Options:} A) True; B) False
\\ \textit{Note:} According to the passage: the west
\item \textbf{Answer:} False
\\ \textit{Options:} A) True; B) False
\\ \textit{Note:} The correct answer is: 1.1 \ensuremath{\times} 1011 metric tonnes
\item \textbf{Answer:} False
\\ \textit{Options:} A) True; B) False
\\ \textit{Note:} The correct answer is: 14 th century
\end{enumerate}

\section*{Long Form Response}
\begin{enumerate}[label=\arabic*.]
\setcounter{enumi}{10}
\item \textbf{Answer:} 115,000
\\ \textit{Note:} Expected answer: 115,000
\item \textbf{Answer:} friendly fire
\\ \textit{Note:} Expected answer: friendly fire
\end{enumerate}

\vfill
\noindent\textit{End of Answer Key}
\end{document}